\subsection{CHEVALLEY 序}

设 \((\mathfrak{g},\mathfrak{h})\) 是 \(\mathbf{Q}\) 上的分裂约化李代数，\(\mathrm{R}\) 是其根系。选取：

\begin{enumerate}
\def\labelenumi{\alph{enumi})}
\item
  一个许可格 \(\mathcal{H}\) 位于 \(\mathfrak{h}\) 中（第 6 号，定义 1）；
\item
  对所有 \(\alpha \in \mathrm{R}\)，取一个格 \(\mathcal{G}^{\alpha}\)，位于 \(\mathfrak{g}^{\alpha}\) 中。
\end{enumerate}

置 \(\mathcal{G} = \mathcal{H} \oplus \sum_{\alpha \in \mathrm{R}} \mathcal{G}^{\alpha}\)。这是 \(\mathfrak{g}\) 中的一个格。记 \(\mathcal{U}\) 为一个 \(\mathbf{Z}\)-子代数，包含于 \(\mathrm{U}(\mathfrak{g})\)，它由 \(\binom{h}{n}\)（\(h \in \mathcal{H}, n \in \mathbf{N}\)）和 \(x^{(n)}\)（\(x \in \mathcal{G}^{\alpha}, \alpha \in \mathrm{R}, n \in \mathbf{N}\)）生成。最后，对 \(\alpha \in \mathrm{R}\) 和 \(x \in \mathfrak{g}^{\alpha} - \{0\}\)，置

\[
w _ {\alpha} (x) = (\exp \operatorname{ad} x) (\exp \operatorname{ad} y) (\exp \operatorname{ad} x),
\]

其中 \(y\) 是 \(\mathfrak{g}^{-\alpha}\) 中满足 \([y,x] = H_{\alpha}\) 的唯一元素。采用这些记号：

定理 2. 以下条件等价：

\begin{enumerate}
\def\labelenumi{(\roman{enumi})}
\item
  存在谢瓦莱系 \((X_{\alpha})_{\alpha \in \mathrm{R}}\)，属于 \((\mathfrak{g},\mathfrak{h})\)，并满足 \(\mathcal{G}_{\alpha} = \mathbf{Z}X_{\alpha}\) 对所有 \(\alpha \in \mathrm{R}\) 成立。
\item
  \(\mathcal{U} \cap \mathfrak{g} = \mathcal{G}\)，并且 \([\mathcal{G}^{\alpha}, \mathcal{G}^{-\alpha}] = \mathbf{Z}H_{\alpha}\) 对所有 \(\alpha \in \mathrm{R}\) 成立。
\item
  对所有 \(\alpha \in \mathrm{R}, x \in \mathcal{G}^{\alpha}, n \in \mathbf{N}\)，自同态 \((\operatorname{ad} x)^n / n!\) 作用于 \(\mathfrak{g}\)，将 \(\mathcal{G}\) 映到 \(\mathcal{G}\)，并且 \([\mathcal{G}^{\alpha}, \mathcal{G}^{-\alpha}] = \mathbf{Z}H_{\alpha}\)。
\item
  对所有 \(\alpha \in \mathrm{R}\) 及每个基向量 \(x\)（属于 \(\mathcal{G}^{\alpha}\)），\(w_{\alpha}(x)\) 将 \(\mathcal{G}\) 映到 \(\mathcal{G}\)（即将 \(\mathcal{G}^{\beta}\) 映到 \(\mathcal{G}^{s_{\alpha}(\beta)}\)，其中 \(\beta \in \mathrm{R}\)）。
\item
  \(\Longrightarrow\) (ii)：令 \((X_{\alpha})_{\alpha \in \mathrm{R}}\) 为 \((\mathfrak{g},\mathfrak{h})\) 中满足 \(\mathcal{G}^{\alpha} = \mathbf{Z}X_{\alpha}\) 对所有 \(\alpha \in \mathrm{R}\) 成立的谢瓦莱系，并令 B 为 R 的一组基。对 \(\alpha \in \mathrm{B}\)，置 \(x_{\alpha} = X_{\alpha},y_{\alpha} = X_{-\alpha}\)。令 \(\mathcal{U}'\) 为第 6 号命题 3 根据 \(\mathcal{H}\)、\(x_{\alpha}\) 和 \(y_{\alpha}\) 关联的双序。显然 \(\mathcal{U}' \subset \mathcal{U}\)。根据命题 3 的推论 3 和 4，都有 \(x^{(n)} \in \mathcal{U}'\)，对所有 \(\alpha \in \mathrm{R}\)、\(x \in \mathcal{G}^{\alpha}\) 及 \(n \in \mathbf{N}\) 成立。因此 \(\mathcal{U} = \mathcal{U}'\)，这证明 (ii)。
\item
  \(\Longrightarrow\) (iii)：这由第 5 号的引理 2 立即得到。
\item
  \(\Longrightarrow\) (iv)：令 \(\alpha \in \mathrm{R}\)，并令 \(x\) 为 \(\mathcal{G}^{\alpha}\) 的一个基向量。由于 \([\mathcal{G}^{\alpha},\mathcal{G}^{-\alpha}] = \mathbf{Z}H_{\alpha}\)，唯一的 \(y \in \mathfrak{g}^{-\alpha}\) 满足 \([y,x] = H_{\alpha}\)，并属于 \(\mathcal{G}^{-\alpha}\)。根据 (iii)，\(\exp \mathrm{ad} x\) 和 \(\exp \mathrm{ad} y\) 保持 \(\mathcal{G}\) 稳定，因而 \(w_{\alpha}(x)\) 也如此。
\item
  \(\Longrightarrow\) (i)：令 B 为 R 的一组基。选取一组基 \(x_{\alpha}\)，属于 \(\mathcal{G}^{\alpha}\)，对所有 \(\alpha \in \mathrm{B}\) 成立。令 \(y_{\alpha} \in \mathcal{G}^{-\alpha}\) 满足 \([y_{\alpha}, x_{\alpha}] = H_{\alpha}\)。根据第 1 节第 5 号的公式 (5)，有 \(y_{\alpha} = w_{\alpha}(x_{\alpha}). x_{\alpha}\)，所以根据 (iv)，\(y_{\alpha}\) 是 \(\mathcal{G}^{-\alpha}\) 的一组基。令 \(\mathcal{G}\) 为 \(\mathfrak{g}\) 中由 \(\mathcal{H}\)、\(x_{\alpha}\) 和 \(y_{\alpha}\) 定义的序（第 6 号，命题 3 的推论 1）。于是 \(\mathcal{G}\) 在 \((\operatorname{ad} x_{\alpha})^{n} / n!\)、\((\operatorname{ad} y_{\alpha})^{n} / n!\) 下稳定（同上），因而也在 \(w_{\alpha}(x_{\alpha})\) 下稳定。
\end{enumerate}

现在令 \(\beta \in \mathrm{R}\)。存在 \(\alpha_0, \alpha_1, \ldots, \alpha_r \in \mathrm{B}\)，使得

\[
\beta = s _ {\alpha_ {r}} s _ {\alpha_ {r - 1}} \dots s _ {\alpha_ {1}} (\alpha_ {0})
\]

（第 VI 章，第 1 节，第 5 号，命题 15）。于是，由 (iv)，\(w_{\alpha_r}(x_{\alpha_r}).w_{\alpha_{r-1}}(x_{\alpha_{r-1}}) \ldots w_{\alpha_1}(x_{\alpha_1})\) 将 \(\mathcal{G}^{\alpha_0}\) 映到 \(\mathcal{G}^\beta\)，并且根据前面的结果，将 \(\mathcal{G} \cap \mathfrak{g}^{\alpha_0}\) 映到 \(\mathcal{G} \cap \mathfrak{g}^\beta\)。由于 \(\mathcal{G} \cap \mathfrak{g}^{\alpha_0} = \mathcal{G}^{\alpha_0}\)（命题 3 (ii)），有 \(\mathcal{G} \cap \mathfrak{g}^\beta = \mathcal{G}^\beta\)。因此

\[
\mathcal {G} = \mathcal {H} \oplus \sum_ {\beta \in \mathrm{R}} (\mathcal {G} \cap \mathfrak {g} ^ {\beta}) = \mathcal {H} \oplus \sum_ {\beta \in \mathrm{R}} \mathcal {G} ^ {\beta} = \mathcal {G}
\]

且命题 3 的推论 4 完成证明。

定义 2. 当定理 2 的条件 (i) 至 (iv) 满足时，称 \(\mathcal{G}\) 为 \((\mathfrak{g}, \mathfrak{h})\) 中的谢瓦莱序。

备注。 \((\mathfrak{g},\mathfrak{h})\) 中的谢瓦莱序总是存在。事实上，谢瓦莱序是形如 \(\mathcal{H}\oplus\sum_{\alpha\in \mathrm{R}}\mathbf{Z}X_{\alpha}\) 的集合，其中 \((X_{\alpha})_{\alpha\in \mathrm{R}}\) 是 \((\mathfrak{g},\mathfrak{h})\) 中的一个谢瓦莱系，而 \(\mathcal{H}\) 是 \(\mathfrak{h}\) 中满足

\[
\mathrm{Q} (\mathrm{R} ^ {\vee}) \subset \mathcal {H} \subset \mathrm{P} (\mathrm{R} ^ {\vee}) \oplus \mathfrak {c}
\]

（其中 c 是 g 的中心）。

定理 3. 保持第 7 号开头的记号，并假设 \(\mathcal{G}\) 是 \((\mathfrak{g},\mathfrak{h})\) 中的谢瓦莱序。

\begin{enumerate}
\def\labelenumi{(\roman{enumi})}
\item
  \(\mathcal{U}\) 是 \(\mathrm{U}(\mathfrak{g})\) 中的 b-序。
\item
  令 B 为 \(\mathrm{R}\) 的一组基，并令 \((X_{\alpha})_{\alpha\in\mathrm{B}\cup(-\mathrm{B})}\) 为 \(\mathfrak{g}\) 中满足 \(\mathcal{G}^{\alpha}=\mathbf{Z}X_{\alpha}\) 对 \(\alpha\in\mathrm{B}\cup(-\mathrm{B})\) 成立的一族元素。\(\mathbf{Z}\)-代数 \(\mathcal{U}\) 由 \(\binom{h}{n}\) 和 \(X_{\alpha}^{(n)}\)（\(h\in\mathcal{H},\alpha\in\mathrm{B}\cup(-\mathrm{B}),n\in\mathbf{N}\)）生成。若 \(\mathfrak{g}\) 是半单的且 \(\mathcal{H}=\mathrm{Q}(\mathrm{R}^{\vee})\)，则 \(\mathbf{Z}\)-代数 \(\mathcal{U}\) 由 \(X_{\alpha}^{(n)}\)（\(\alpha\in\mathrm{B}\cup(-\mathrm{B}),n\in\mathbf{N}\)）生成。
\item
  令 B 为 R 的一组基，\(R_{+}\) 为相应的正根集合，\(R_{-} = -R_{+}\)，\(n_{+} = \sum_{\alpha \in R_{+}} g^{\alpha}\)，\(n_{-} = \sum_{\alpha \in R_{-}} g^{\alpha}\)。则
\end{enumerate}

\[
\mathcal {U} = (\mathcal {U} \cap \mathrm{U} (\mathfrak {n} _ {-})). (\mathcal {U} \cap \mathrm{U} (\mathfrak {h})). (\mathcal {U} \cap \mathrm{U} (\mathfrak {n} _ {+})).
\]

令 \((h_i)_{i\in \mathrm{I}}\) 为 \(\mathcal{H}\) 的一组基。对所有 \(\alpha \in \mathrm{R}\)，令 \(X_{\alpha}\) 为 \(\mathcal{G}^{\alpha}\) 的一组基。给集合 \(\mathrm{I}\cup \mathrm{R}\) 赋予一个全序（我们假定 \(\mathrm{I}\cap \mathrm{R} = \emptyset\)）。对 \(\lambda \in \mathrm{I}\cup \mathrm{R}\) 及 \(n\in \mathbf{N}\)，置 \(e_{\lambda}^{\langle n\rangle} = \binom{h_{\lambda}}{n}\)（若 \(\lambda \in \mathrm{I}\)），置 \(e_{\lambda}^{\langle n\rangle} = X_{\lambda}^{(n)}\)（若 \(\lambda \in \mathrm{R}\)）。则乘积 \(\prod_{\lambda \in \mathrm{I}\cup \mathrm{R}}e_{\lambda}^{\langle n_{\lambda}\rangle}\) 当 \((n_{\lambda})\) 属于 \(\mathbf{N}^{\mathrm{I}\cup \mathrm{R}}\) 时，构成 \(\mathbf{Z}\)-模 \(\mathcal{U}\) 的一组基。乘积 \(\prod_{\lambda \in \mathrm{I}}\binom{h_{\lambda}}{n_{\lambda}}\) 当 \((n_{\lambda})\) 属于 \(\mathbf{N}^{\mathrm{I}}\) 时，构成 \(\mathbf{Z}\)-模 \(\mathcal{U}\cap \mathrm{U}(\mathfrak{h})\) 的一组基。乘积 \(\prod_{\lambda \in \mathrm{R}_{+}}X_{\lambda}^{(n_{\lambda})}\) 当 \((n_{\lambda})\) 属于 \(\mathbf{N}^{\mathrm{R}_{+}}\) 时，构成 \(\mathbf{Z}\)-模 \(\mathcal{U}\cap \mathrm{U}(\mathfrak{n}_{+})\) 的一组基。

令 B 和 \((X_{\alpha})_{\alpha\in\mathrm{B}\cup(-\mathrm{B})}\) 如 (ii) 中所述，并满足 \([X_{-\alpha},X_{\alpha}]=H_{\alpha}\)。令 \(\mathcal{U}'\) 为一个 \(\mathbf{Z}\)-子代数，包含于 \(\mathrm{U}(\mathfrak{g})\)，它由 \(\binom{h}{n}\) 和 \(X_{\alpha}^{(n)}\) \((h\in\mathcal{H},\alpha\in\mathrm{B}\cup(-\mathrm{B}),n\in\mathbf{N})\) 生成。在定理 2 的证明中，从 (i) \(\Longrightarrow\) (ii) 的部分，我们已经看到 \(\mathcal{U}'\) 等于 \(\mathcal{U}\)，并且是 \(\mathrm{U}(\mathfrak{g})\) 中的双序。这证明 (i) 及 (ii) 的第一个断言；第二个断言由引理 4 (ii) 得到。断言 (iii) 由定理 1（第 3 号）和命题 3（第 6 号）得到。

